\begin{partdivider}{II}{LE GRAND BASCULEMENT}\end{partdivider}

\bookchap{4}{Lorsque les machines commencent à penser avec nous}

Pendant des décennies, notre rapport aux ordinateurs a obéi à une règle simple et asymétrique : nous posions une question, ou nous donnions un ordre, et la machine exécutait. Si nous saisissions une formule dans un tableur, le logiciel calculait. Si nous entrions une requête dans un moteur de recherche, il nous renvoyait une liste de pages existantes. La machine était un exécutant prodigieux, rapide et docile, mais dépourvu d'initiative. Elle ne proposait rien ; elle répondait.

Ce schéma a commencé à vaciller avec les progrès de l'apprentissage automatique (le \emph{machine learning}). Plutôt que de programmer explicitement une machine pour qu'elle reconnaisse un chat sur une image — une tâche impossible à réduire à des règles fixes —, nous lui avons montré des milliers de photos de chats, en lui laissant découvrir elle-même les régularités. Le fonctionnement de la machine est devenu moins transparent pour nous, mais plus efficace. La machine classait, triait, prédisait. Elle restait cependant dans un schéma de restitution : à partir de données passées, elle inférait une étiquette ou une probabilité.

Pour reprendre une analogie simple : l'IA traditionnelle fonctionnait comme un archiviste méticuleux. Vous lui demandiez un dossier spécifique selon des critères précis, et il vous le trouvait immanquablement dans les étagères, ou classait les nouveaux documents dans les bonnes cases. L'IA générative, elle, ressemble à un stagiaire créatif et polyglotte : vous lui donnez un ensemble de notes éparses et un objectif de haut niveau, et il rédige pour vous une synthèse inédite, en proposant un angle ou une formulation que vous n'aviez pas envisagés. Le saut n'est pas de la recherche à la vitesse, mais de la restitution à la proposition.

C'est un fait désormais visible : avec l'apparition des réseaux de neurones profonds (le \emph{deep learning}), puis des modèles génératifs à très grande échelle (les LLM, ou grands modèles de langage), une rupture s'est produite dans la nature des sorties que la machine est capable de produire. La machine ne se contente plus de classer ou de prédire : elle génère.

Cette évolution technique, même si elle est encore imparfaite, modifie profondément notre rapport à l'outil. Car un système qui génère du langage, qui peut reformuler une idée, l'articuler dans un style différent, ou en tirer les implications logiques, ne fonctionne plus comme un dictionnaire ou une base de données. Il fonctionne, dans une certaine mesure, comme un interlocuteur.

Il faut ici prendre la mesure d'un changement conceptuel essentiel. Jusqu'à présent, nous adaptions notre pensée au langage de la machine. Nous apprenions à coder, à formuler des requêtes booléennes, à utiliser des interfaces complexes. Avec les modèles génératifs, c'est la machine qui s'adapte au nôtre. Elle comprend — ou du moins traite — le langage naturel. Nous pouvons lui parler comme nous parlerions à un collaborateur, en lui demandant de «résumer ce texte en mettant l'accent sur les risques juridiques» ou de «proposer trois arguments contre cette thèse».

On observe que le moteur de recherche nous renvoyait vers des documents existants, produits par des humains. L'IA générative, elle, synthétise une réponse nouvelle à partir de l'ensemble des données sur lesquelles elle a été entraînée. Elle ne nous dit pas où trouver l'information ; elle produit une formulation qui y répond, fût-ce parfois avec des erreurs ou des approximations.

C'est cette capacité de synthèse et de formulation qui constitue la véritable rupture. Non pas parce que la machine «comprend» au sens philosophique du terme — une question sur laquelle les chercheurs débattent encore —, mais parce que ses productions sont, pour la première fois, suffisamment cohérentes, contextualisées et adaptées à notre intention pour que nous puissions \emph{penser avec elles}.

Prenons une analogie simple. Lorsqu'un écrivain utilise un dictionnaire, il cherche un mot qu'il connaît déjà ou dont il vérifie l'orthographe. Le dictionnaire est un outil de vérification. Lorsqu'un écrivain s'entretient avec un confrère pour surmonter un blocage, il engage un dialogue : il propose une idée, le confrère lui répond, l'écrivain réajuste sa pensée, et de cette interaction émerge une formulation qui n'existait pas avant. L'IA générative, dans son usage le plus fécond, se rapproche davantage de la seconde situation que de la première. Elle ne se contente pas de vérifier ; elle participe, à sa manière, à l'élaboration.

Il est crucial de ne pas surévaluer cette participation. La machine n'a ni intention propre, ni conscience, ni compréhension du monde au sens humain. Ses réponses sont le produit de calculs statistiques extrêmement sophistiqués sur des probabilités de mots. Mais du point de vue de l'utilisateur qui cherche à avancer dans un raisonnement, la distinction pratique importe parfois moins que l'effet produit : la machine lui permet d'avancer plus vite, de contourner un obstacle, d'envisager une perspective qu'il n'avait pas considérée.

Ce basculement explique pourquoi l'IA générative a été adoptée si rapidement, bien au-delà de la sphère technique. Elle a franchi la frontière de l'outil spécialisé pour devenir un système avec lequel on peut dialoguer, et non plus simplement commander. Cependant, dès lors que la machine ne se contente plus d'exécuter mais propose, une nouvelle question se pose : que se passe-t-il lorsque l'outil commence à prendre des initiatives ? C'est ce glissement, de l'outil vers l'agent, que nous devons maintenant explorer.

\bookchap{5}{La fin de l'ordinateur comme simple outil}

Le vocabulaire que nous employons pour parler de l'IA en dit long sur la manière dont notre relation à la machine est en train de se transformer. Pendant longtemps, nous parlions d'«outils informatiques» ou de «logiciels». Aujourd'hui, le langage s'est déplacé : nous parlons d'«assistants», de «copilotes» ou d'«agents».

Ce glissement sémantique n'est pas anodin. Il traduit un changement fonctionnel majeur, que nous pouvons décomposer en plusieurs étapes.

La première étape est celle de \textbf{l'outil}. Un traitement de texte, un tableur ou un logiciel de dessin sont des outils. Ils sont puissants, mais passifs. Ils ne font rien tant que nous ne leur avons pas donné une instruction explicite. Leur valeur réside dans leur capacité à exécuter avec précision et rapidité ce que nous avons défini.

La deuxième étape est celle de \textbf{l'assistant}. Il ne se contente plus d'exécuter, il peut réagir à un contexte. Un correcteur orthographique qui souligne une erreur en temps réel est un assistant rudimentaire. Il anticipe un besoin (la correction) sans que nous ayons à le solliciter explicitement pour chaque mot. Il nous facilite la tâche en arrière-plan.

La troisième étape, celle que nous venons d'entrapercevoir avec les modèles génératifs, est celle du \textbf{copilote}. Le copilote ne se contente plus de corriger, il propose. Si un développeur tape le début d'une fonction, le copilote peut en deviner la suite et lui soumettre un bloc de code complet. Si un rédacteur écrit l'amorce d'un paragraphe, le copilote peut lui proposer une suite logique. Le copilote ne remplace pas l'humain, il travaille

\emph{avec} lui, dans le flux de son activité. Il a une forme d'initiative, mais encadrée : il propose, l'humain décide d'accepter, de modifier ou de rejeter.

La quatrième étape, qui se dessine actuellement, est celle de \textbf{l'agent}. L'agent va plus loin que le copilote. On ne lui demande plus seulement de proposer la ligne suivante, on lui confie un objectif. Par exemple : «Organise mon voyage à Londres pour la semaine prochaine, réserve les hôtels dans ce budget, et synthétise les avis sur ces restaurants.» L'agent va décomposer cette requête en sous-tâches, chercher l'information, comparer les options, effectuer les réservations et présenter un résumé. Il planifie, il exécute, il rend compte.

On observe que nous passons progressivement d'un modèle où l'humain pilotait la machine pas à pas, à un modèle où l'humain fixe une intention, et la machine conçoit et exécute un plan d'action pour y répondre.

Cette évolution est séduisante en termes de productivité. Elle est aussi profondément déstabilisante. Car déléguer une intention, et non plus une instruction, change la nature du contrôle que nous exerçons.

Lorsque je donne une instruction précise à un outil, je sais exactement ce qu'il va faire. Lorsque je confie une intention à un agent, je ne maîtrise plus toutes les étapes intermédiaires. L'agent a fait des choix, il a interprété ma demande, il a pris des raccourcis ou des initiatives que je n'ai pas explicitement validés. Cela soulève des questions de confiance, de supervision et de responsabilité.

Prenons un exemple concret. Un dirigeant demande à un agent IA de «préparer un rapport sur les risques concurrentiels du secteur, avec des recommandations stratégiques». L'agent rédige un rapport cohérent, citant des études de marché, identifiant des tendances et proposant des axes d'action. Le rapport est plausible, bien structuré et rapide à produire. Mais le dirigeant n'a pas vérifié les sources. Il n'a pas supervisé le raisonnement. Il a délégué non seulement l'exécution, mais une partie de l'analyse.

\begin{aiquote}
Que se passe-t-il lorsque nous déléguons le raisonnement intermédiaire à une machine ?
\end{aiquote}

Nous gagnons du temps. Nous accédons peut-être à des analyses que nous n'aurions pas eu le temps de faire. Mais nous risquons aussi de perdre en compréhension. Le rapport devient une boîte noire : le résultat est là, mais le cheminement pour y parvenir m'échappe en partie.

C'est la différence entre conduire une voiture et être conduit. Lorsque je conduis, je décide de chaque virage, je perçois la route, je comprends pourquoi j'ai choisi cet itinéraire. Lorsque je suis conduit, j'arrive à destination — peut-être plus vite et plus sûrement — mais si le chauffeur se trompe, je ne m'en apercevrai peut-être que trop tard, parce que je n'ai pas suivi le trajet.

L'hypothèse est la suivante : le passage de l'outil à l'agent ne modifie pas seulement notre productivité, il modifie notre rapport à l'action et à la compréhension. Plus nous déléguons l'initiative, plus nous risquons de nous trouver en situation de dépendance cognitive à l'égard de systèmes dont nous ne maîtrisons pas entièrement les inférences.

Ce risque ne signifie pas qu'il faille renoncer aux agents. Les humains délèguent depuis toujours des tâches complexes à d'autres humains (collaborateurs, conseillers, experts) sans en maîtriser tous les détails. Mais avec les humains, nous avons développé des mécanismes de confiance, de reddition de comptes et de jugement mutuel qui se sont affinés au fil des millénaires. Avec les agents IA, ces mécanismes restent à inventer, à adapter ou à réinventer.

La question de la délégation est d'autant plus cruciale que cette évolution technologique s'accompagne d'une autre transformation, économique cette fois : l'effondrement potentiel du coût de production.

\bookchap{6}{Le jour où produire devient presque gratuit}

L'une des conséquences les plus visibles de l'IA générative est la baisse brutale du temps et du coût nécessaires pour produire un grand nombre de biens cognitifs et créatifs.

Pendant des siècles, produire un texte, un logiciel, une image, une analyse, une traduction ou un cours exigeait du temps humain. Ce temps avait un coût. Ce coût structurait l'économie de la production intellectuelle. Un article de presse coûtait cher à produire parce qu'il fallait un journaliste pour le rédiger. Un logiciel coûtait cher parce qu'il fallait des développeurs pour l'écrire ligne par ligne. Une illustration coûtait cher parce qu'elle exigeait le talent et les heures de travail d'un artiste.

C'est un fait : avec l'IA, le coût marginal de production de certaines de ces réalisations tend vers une valeur très faible, parfois proche de zéro. Il est aujourd'hui possible de générer en quelques secondes et pour quelques centimes une image en haute définition, une traduction fluide, un résumé structuré ou un programme fonctionnel.

Il faut prendre la mesure de ce que cette évolution signifie, non pas en termes technologiques, mais en termes économiques et sociaux. Pendant toute l'histoire de l'humanité, la rareté des biens et des services a été le postulat de base de l'organisation économique. Les systèmes de prix, les salaires, les marchés se sont construits pour allouer des ressources rares. La production intellectuelle n'échappait pas à cette règle. Avec l'IA, ce postulat vacille.

L'hypothèse est la suivante : si la production de certaines réalisations intellectuelles ou créatives devient quasi gratuite, l'un des piliers sur lesquels repose la valorisation de ces activités se trouve ébranlé.

Considérons quelques exemples concrets.

Premier exemple : la rédaction de contenu. Une petite entreprise qui souhaitait produire régulièrement des articles de blog, des fiches produits ou des newsletters devait soit embaucher un rédacteur, soit sous-traiter à une agence, soit y consacrer son propre temps. Aujourd'hui, elle peut générer une dizaine de variantes en quelques minutes avec une IA, pour un coût dérisoire. Le contenu n'a pas disparu, il est même devenu abondant.

Deuxième exemple : la programmation. Il y a encore peu de temps, développer une application simple nécessitait l'intervention d'un développeur, dont le temps était facturé cher. Avec les copilotes de programmation, un individu maîtrisant les bases du code peut multiplier sa vitesse de production par cinq ou par dix. Le coût de développement d'un prototype s'effondre.

Troisième exemple : la traduction. Une entreprise qui devait traduire son site en cinq langues engageait autrefois des traducteurs professionnels. Aujourd'hui, une traduction automatique post-éditée par l'IA permet de réaliser l'essentiel du travail à une fraction du coût.

Dans chacun de ces cas, la production n'a pas été supprimée, elle a été considérablement accélérée et rendue moins coûteuse. L'abondance de contenu, de code et de traductions est en train d'exploser.

Mais cette explosion soulève une question économique fondamentale :

\begin{aiquote}
Que devient la valeur lorsque produire devient beaucoup plus facile ?
\end{aiquote}

Il serait naïf de croire que la valeur disparaît. Elle se déplace. Et ce déplacement est prévisible, même si ses conséquences exactes restent à observer.

Premièrement, \textbf{la valeur se déplace de la production vers la curation et le jugement.} Lorsque le contenu est rare, le simple fait de produire a de la valeur. Lorsque le contenu est abondant, la capacité à trier, à filtrer, à évaluer ce qui est bon, pertinent et fiable prend le dessus. Dans un monde noyé sous les textes et les images générés, celui qui sait dire lequel mérite attention détient une nouvelle forme de valeur. Deuxièmement, \textbf{la valeur se déplace vers l'originalité, la singularité et la profondeur.} Si les analyses standards, les résumés ou les codes génériques deviennent des commodités, ce qui sort du lot — une perspective véritablement nouvelle, un raisonnement qui ne se contente pas de synthétiser le sens commun, une créativité qui ne se réduit pas à la recombinaison statistique — pourrait devenir plus précieux que jamais. Le coût de l'ordinaire baisse ; la valeur de l'extraordinaire, ou du moins du singulier, pourrait augmenter.

Troisièmement, \textbf{la valeur se déplace vers la relation et la confiance.} Pourquoi paie-t-on encore pour un coach sportif alors qu'on pourrait trouver tous les exercices gratuits sur Internet ? Pourquoi consulte-t-on un thérapeute alors que des livres expliquent parfaitement les mécanismes psychologiques ? Parce que la valeur ne réside plus dans l'information elle-même, mais dans la présence, l'adaptation à la personne, la relation de confiance et la responsabilisation. De même, un avocat, un consultant ou un enseignant pourraient voir la part «informationnelle» de leur métier s'effondrer, tandis que la part relationnelle, stratégique et humaine prendrait le relais.

Notre interprétation est la suivante : l'effondrement du coût de production ne signe pas la fin du travail intellectuel. Il modifie la répartition de la valeur au sein de ce travail. Les tâches d'exécution cognitives — formuler, coder, calculer, traduire — perdent de leur valeur économique. Les tâches de supervision, de jugement, de relation et de définition des objectifs en gagnent potentiellement.

Il faut cependant faire face à une réalité plus brute. Ce déplacement de la valeur sera probablement vécu, pendant une période transitoire, comme une destruction brutale de revenus pour ceux dont l'activité reposait précisément sur l'exécution cognitive. Un traducteur professionnel, un illustrateur, un rédacteur technique ou un développeur junior dont le travail est principalement d'exécution verront leurs marchés se contracter ou se transformer très rapidement. Le déplacement de la valeur vers le jugement ou la relation ne créera pas mécaniquement des emplois pour ceux qui viennent d'en perdre. La transition peut être douloureuse et inégalitaire.

De plus, l'abondance de production à faible coût pose un problème de nature différente : celui de la confiance. Si n'importe qui peut générer en une heure un rapport de cinquante pages qui a l'air autorisé, n'importe qui peut aussi inonder l'espace public de textes fallacieux, d'images trompeuses ou de vidéos contrefaites. L'abondance de production rend plus aiguë la nécessité de vérifier, d'auditer et d'authentifier.

Nous rejoignons ici l'une des tensions centrales de ce livre :

\begin{eplist}
\item D'un côté, l'IA démocratise l'accès à la production intellectuelle. Elle permet à des individus et de petites organisations de réaliser ce qui était autrefois réservé aux grandes équipes ou aux experts bien dotés.
\item De l'autre, elle noie l'espace commun sous une production de masse, rendant plus difficile la distinction entre le vrai et le faux, le pertinent et le généré, le fiable et le simulé.
\end{eplist}
Le jour où produire devient presque gratuit n'est donc pas le jour où tout devient simple. C'est, au contraire, peut-être le jour où la question de la valeur, de la confiance et du jugement devient plus aiguë que jamais.

Ce basculement de la production et de la valeur nous oblige à regarder en face la question qui inquiète et fascine le plus : celle du travail. Car si l'IA transforme la manière dont nous produisons, elle transforme aussi, nécessairement, la manière dont nous travaillons. C'est ce que nous explorerons dans la troisième partie de ce livre.
